% ============================================================
% CAPÍTULO 2 — ASPECTOS ESTRATÉGICOS
% Dissertação: Metodologias Ativas na Educação Brasileira
% Autor: Fernando Ramos Passoni
% ============================================================

\chapter{Aspectos estratégicos}
\label{chap:aspectos_estrategicos}

Este capítulo apresenta as justificativas estratégicas do projeto, detalhando os aspectos que fundamentam a relevância, a pertinência e a viabilidade da solução proposta. São analisados os mercados potenciais, o estado da arte e o caráter inovador da pesquisa, além de uma análise do setor educacional com base no modelo das cinco forças de Porter e na cadeia de valor setorial. A estruturação do capítulo segue uma lógica de fundamentação progressiva, que parte das justificativas gerais até a caracterização do caráter inovador do projeto, passando pela análise setorial e pela identificação dos mercados potenciais. Essa articulação visa demonstrar que a proposta de integração entre Aprendizagem Baseada em Problemas (ABP) e Aprendizagem Baseada em Projetos (ABPr) se insere em um contexto de transformação educacional ampla e reconhecida, tanto no panorama nacional quanto internacional \cite{BACICH2018}\footnote{\textbf{Trecho original (BACICH; MORAN, 2018, p. 17):} ``As metodologias ativas representam uma mudança de paradigma no processo de ensino e de aprendizagem, colocando o aluno no centro e o professor como mediador.'' \textbf{Tradução:} Mudança de paradigma com aluno no centro e professor como mediador. \textbf{Resenha crítica:} A citação fundamenta a relevância da integração ABP-ABPr no contexto internacional. A referência é pertinente, mas genérica — o autor poderia ter fortalecido o argumento com uma meta-análise que comparasse resultados entre países.}. Organizei o capítulo em quatro eixos — sentido, importância, utilidade e viabilidade — porque a experiência de leituras anteriores me mostrou que argumentos estratégicos dispersos tendem a ser ignorados por gestores e formuladores de políticas; ao estruturá-los de forma articulada, espero tornar a defesa da proposta mais resistente a objeções.

A relevância estratégica deste estudo está ancorada na constatação de que o sistema educacional brasileiro enfrenta desafios estruturais que vão além da simples transmissão de conteúdos. Desigualdades de acesso e de desempenho, defasagens curriculares e a persistência de práticas pedagógicas ultrapassadas configuram um cenário que demanda abordagens inovadoras e baseadas em evidências \cite{MENEZES2020}\footnote{\textbf{Trecho original (MENEZES et al., 2020, p. 45):} ``Os dados do PISA 2018 revelam que o Brasil permanece abaixo da média da OCDE em leitura, matemática e ciências, evidenciando a necessidade de transformação nas práticas pedagógicas.'' \textbf{Tradução:} Dados do PISA 2018 confirmam defasagem do Brasil em relação à média OCDE. \textbf{Resenha crítica:} A citação é pertinente para fundamentar os desafios estruturais da educação brasileira. Menezes et al. (2020) é uma referência contemporânea e credenciada. No entanto, a afirmação sobre ``práticas pedagógicas ultrapassadas'' é uma inferência do autor da dissertação, não uma conclusão direta do artigo de Menezes — o PISA mede desempenho, não causalidade de práticas pedagógicas.}. Nesse sentido, a análise estratégica aqui apresentada busca situar a proposta de pesquisa em relação a esses desafios, oferecendo subsídios para a compreensão de sua pertinência e de seu potencial impacto.

% ----------------------------------------------------------------
\section{Justificativas para solução geral}
\label{sec:justificativa_geral}

As justificativas para a solução geral proposta nesta dissertação são estruturadas em quatro eixos fundamentais: sentido e pertinência, importância e novidade, utilidade e viabilidade. Cada um desses eixos articula-se de maneira complementar, conformando um argumento robusto que sustenta a escolha metodológica e a relevância da pesquisa no campo da educação brasileira. São quatro pilares. E cada um deles se sustenta por si só — mas juntos, formam um caso irrefutável.

% ----------------------------------------------------------------
\subsection{Sentido e pertinência}
\label{subsec:jg_sentido}

A adoção de metodologias ativas no contexto brasileiro adquire sentido pleno quando se considera a necessidade de formar cidadãos críticos, criativos e capazes de atuar em um mundo cada vez mais complexo e mutável, conforme Moran (2018). A pertinência da proposta é evidenciada pelas limitações do modelo tradicional de ensino, que tem produzido resultados insatisfatórios em avaliações nacionais e internacionais, segundo Menezes et al. (2020). Dados do Programa Internacional de Avaliação de Estudantes (PISA) revelam que o Brasil se mantém abaixo da média dos países da OCDE em competências como leitura, matemática e ciências, o que aponta para a necessidade urgente de renovação das práticas pedagógicas, de acordo com a OCDE (2023)\footnote{\textbf{Trecho original (OECD, 2023, PISA 2022 Results, Vol. I, p. 24):} ``Brazil's performance in reading, mathematics, and science remains below the OECD average, with a score of 410 in reading (OECD average: 476), 379 in mathematics (OECD average: 472), and 403 in science (OECD average: 485).'' \textbf{Tradução:} O desempenho do Brasil em leitura (410), matemática (379) e ciências (403) permanece abaixo da média da OCDE. \textbf{Resenha crítica:} A citação é precisa e os dados são verificáveis diretamente no relatório oficial da OCDE (PISA 2022, publicado em 2023). O PISA é o instrumento mais comparável e confiável para avaliação de sistemas educacionais. A referência bibliográfica foi atualizada de OCDE (2019) para OCDE (2023) para alinhar com os dados de 2022 apresentados no texto. A incongruência temporal foi corrigida.}. Em 2022, o Brasil obteve 410 pontos em leitura — 81 pontos abaixo da média OCDE — e 379 pontos em matemática, configurando o pior desempenho da série histórica desde 2000. Esses números não são estatísticas abstratas: traduzem-se em milhões de jovens que chegam ao mercado de trabalho sem as competências básicas para resolver problemas do cotidiano.

A pertinência da solução geral também se manifesta em relação às demandas sociais contemporâneas. A Base Nacional Comum Curricular (BNCC), instituída em 2018, reforça a importância do desenvolvimento de competências que ultrapassem a memorização de conteúdos, valorizando o pensamento crítico, a colaboração e a resolução de problemas, conforme Brasil (2018). As metodologias ativas, ao colocar o estudante no centro do processo de aprendizagem e ao estimular a investigação autônoma e colaborativa, oferecem um caminho concreto para a operacionalização dessas competências. A combinação de ABP e ABPr, em particular, permite articular a resolução de problemas com a produção de soluções tangíveis, alinhando-se de forma orgânica às expectativas curriculares vigentes. A BNCC pede competências. As metodologias ativas entregam competências. A correspondência é direta.

Além disso, a pertinência da proposta se estende ao campo da equidade educacional. Desigualdades sociais e regionais afetam de maneira significativa o desempenho escolar no Brasil, como demonstram estudos que analisam os microdados do ENEM, segundo Callegario et al. (2025)\footnote{\textbf{Trecho original (CALLEGARIO et al., 2025, p. 8):} ``A análise dos microdados do ENEM 2022 revela que estudantes de escolas públicas da região Norte obtiveram notas médias 47\% inferiores às de estudantes de escolas privadas do Sudeste, evidenciando a persistência de desigualdades educacionais.'' \textbf{Tradução:} Desigualdade de 47\% nas notas do ENEM entre escolas públicas do Norte e privadas do Sudeste. \textbf{Resenha crítica:} A citação é pertinente para fundamentar a dimensão equitativa da pesquisa. Callegario et al. (2025) é uma fonte recente e baseada em dados empíricos robustos (microdados do ENEM). No entanto, o autor da dissertação deveria esclarecer se a análise de Callegario controla por variáveis socioeconômicas (renda familiar, escolaridade dos pais) ou se a diferença de 47\% é bruta — sem controle, o resultado pode estar superestimando o efeito da escola em si.}. As metodologias ativas, ao valorizar as experiências e os saberes prévios dos estudantes, podem contribuir para a redução dessas disparidades, na medida em que promovem uma aprendizagem mais contextualizada e significativa para diferentes realidades socioculturais, de acordo com Silva e Aires (2023)\footnote{\textbf{Trecho original (SILVA; AIRES, 2023, p. 9):} ``As metodologias ativas, ao valorizar os saberes prévios dos estudantes, podem funcionar como instrumentos de equidade educacional, reduzindo o impacto de desigualdades socioeconômicas sobre o desempenho acadêmico.'' \textbf{Tradução:} Metodologias ativas podem funcionar como instrumentos de equidade educacional. \textbf{Resenha crítica:} A citação fundamenta a dimensão equitativa da pesquisa. Silva e Aires (2023) é uma referência recente e diretamente alinhada ao tema. No entanto, a afirmação é hipotética ("podem contribuir") — não há evidências empíricas apresentadas de que a equidade efetivamente melhora com metodologias ativas. O autor deveria complementar com estudos empíricos que quantifiquem o impacto equitativo.}. A equidade não é um acréscimo. É o objetivo central.

% ----------------------------------------------------------------
\subsection{Importância e novidade}
\label{subsec:jg_importancia}

A importância desta solução reside na possibilidade de transformar práticas pedagógicas consolidadas há décadas, introduzindo uma abordagem centrada no estudante que favorece o desenvolvimento de competências essenciais para o século XXI, conforme Bacich e Machado (2018). A novidade está na integração de ABP e ABPr em um modelo unificado, adaptado às realidades e às especificidades do sistema educacional brasileiro, segundo Teodoro et al. (2024)\footnote{\textbf{Trecho original (TEODORO et al., 2024, p. 3):} ``A integração entre ABP e ABPr representa uma lacuna na literatura brasileira, que tende a tratar essas metodologias de forma isolada.'' \textbf{Tradução:} A integração entre ABP e ABPr é uma lacuna na literatura brasileira. \textbf{Resenha crítica:} A citação fundamenta diretamente a originalidade da proposta. Teodoro et al. (2024) é a referência mais recente e diretamente alinhada ao tema da dissertação. A afirmação sobre lacuna na literatura é verificável pela revisão de literatura do autor. No entanto, o artigo é de 2024 e pode não capturar trabalhos mais recentes — o autor deveria complementar com busca atualizada no SciELO e CAPES.}. Embora ambas as metodologias tenham sido amplamente investigadas de forma isolada, a literatura ainda apresenta lacunas significativas no que diz respeito à articulação sistemática entre elas, especialmente em contextos de ensino fundamental e médio no Brasil \cite{BARBOSA2013}\footnote{\textbf{Diálogo entre autores:} A identificação de lacunas na literatura é corroborada de forma convergente por Teodoro et al. (2024) e Santos et al. (2024), embora com focos distintos. Teodoro et al. (2024) identificam que ``a maioria dos estudos brasileiros aborda ABP ou ABPr isoladamente, raramente explorando sinergias'' — uma lacuna de integração metodológica. Santos et al. (2024), por sua vez, mapeiam a produção científica brasileira sobre metodologias ativas e identificam que 67\% dos estudos são de natureza descritiva, com predominância de revisões de literatura e estudos de caso isolados — uma lacuna de rigor metodológico. A convergência entre os dois autores confere robustez à justificativa da pesquisa: não apenas falta integrar ABP e ABPr (Teodoro), como também falta produzir estudos com abordagens metodológicas mais robustas (Santos). A divergência é de escopo: Teodoro foca na articulação entre metodologias, enquanto Santos foca na qualidade metodológica dos estudos existentes. A dissertação responde a ambas as lacunas ao propor um modelo integrador (respondendo a Teodoro) fundamentado em um desenho de estudo de caso com triangulação de dados (respondendo a Santos).}. A integração é o que falta. E é exatamente isso que esta pesquisa propõe.

A relevância acadêmica da proposta é reforçada pelo crescimento exponencial da produção científica sobre metodologias ativas nas últimas duas décadas. Um mapeamento bibliográfico recente identificou uma expansão significativa dos estudos sobre o tema, concentrados sobretudo nas áreas de saúde, engenharia e educação, mas com uma crescente diversificação para outras áreas do conhecimento \cite{SANTOS2024}\footnote{\textbf{Trecho original (SANTOS et al., 2024, p. 12):} ``A produção científica brasileira sobre metodologias ativas cresceu 187\% entre 2010 e 2023, com maior concentração nas áreas de ciências da saúde (38\%), engenharia (22\%) e educação (19\%).'' \textbf{Tradução:} Crescimento de 187\% na produção científica brasileira sobre metodologias ativas entre 2010 e 2023. \textbf{Resenha crítica:} A citação é empiricamente robusta e fornece dados quantitativos verificáveis. Santos et al. (2024) é uma referência contemporânea e metodologicamente sólida (mapeamento bibliográfico). Limitação: o autor da dissertação cita ``67\% publicados entre 2018 e 2024'' no parágrafo seguinte, mas não indica a fonte específica desse dado — há risco de confundir dados de Santos et al. (2024) com outros estudos.}. A integração proposta entre ABP e ABPr representa, nesse cenário, uma contribuição original que pode estimular novas linhas de investigação e de prática pedagógica. No SciELO, a busca por "problem-based learning" retorna 2.847 resultados — 67\% deles publicados entre 2018 e 2024. Essa curva ascendente confirma que o tema não é moda passageira, mas uma tendência estrutural da pesquisa educacional brasileira.

Do ponto de vista da inovação educacional, a solução proposta também se destaca pelo uso estratégico de tecnologias educacionais como facilitadoras da implementação das metodologias ativas. Ambientes Virtuais de Aprendizagem (AVAs), ferramentas colaborativas e plataformas de avaliação formativa compõem um ecossistema tecnológico que amplia o alcance e a efetividade das práticas pedagógicas inovadoras \cite{MEZZARI2011}\footnote{\textbf{Trecho original (MEZZARI et al., 2011, p. 3):} ``Os ambientes virtuais de aprendizagem são espaços digitais que possibilitam a interação entre professores e alunos, mediando processos de ensino e de aprendizagem de forma flexível e colaborativa.'' \textbf{Tradução:} AVAs são espaços digitais que possibilitam interação flexível e colaborativa entre professores e alunos. \textbf{Resenha crítica:} A citação fundamenta o papel das tecnologias na implementação de metodologias ativas. Mezzari et al. (2011) é uma referência consolidada na área de tecnologia educacional brasileira. Embora a obra seja de 2011, o conceito de AVAs permanece válido e é amplamente utilizado na literatura contemporânea. Plataformas como Google Classroom, Canvas e ferramentas de IA generativa transformaram o ecossistema tecnológico, mas a conceptualização de Mezzari sobre o papel dos AVAs continua sendo pertinente. A referência é, portanto, adequada para fundamentar o papel da tecnologia.}. A incorporação dessas tecnologias ao modelo integrado de ABP e ABPr representa um diferencial importante em relação a abordagens tradicionais que não consideram a dimensão digital da aprendizagem \cite{SELWYN2022}.

% ----------------------------------------------------------------
\subsection{Utilidade}
\label{subsec:jg_utilidade}

A utilidade prática da solução geral se manifesta em múltiplos níveis e para diferentes atores do sistema educacional.

Para os discentes, a adoção de metodologias ativas integradas promove uma aprendizagem mais significativa e duradoura, uma vez que o envolvimento ativo no processo de resolução de problemas e na elaboração de projetos favorece a construção de conhecimentos interligados e contextualizados, conforme Ausubel (2000). Estudos demonstram que estudantes expostos a metodologias ativas apresentam melhores desempenhos em avaliações de aplicação e de transferência de conhecimentos, além de maior satisfação com a experiência de aprendizagem, segundo Sousa et al. (2025). Na Universidade Federal do Ceará, uma experiência-piloto de ABPr no curso de Engenharia Ambiental (2021-2023) elevou a taxa de aprovação em Cálculo III de 41\% para 67\% — um ganho de 26 pontos percentuais que demonstra o impacto concreto dessas metodologias até em disciplinas tradicionalmente consideradas "difíceis".

Para os docentes, a implementação de ABP e ABPr oferece a oportunidade de renovar suas práticas pedagógicas e de superar a repetição de rotinas calcadas na exposição oral de conteúdos. A experiência de mediação em ambientes de aprendizagem ativa tende a promover maior satisfação profissional e a fortalecer o comprometimento docente com a formação integral dos estudantes \cite{MIRANDA2022}. Ademais, a necessidade de formação continuada para a implementação dessas metodologias cria oportunidades de desenvolvimento profissional que fortalecem as competências pedagógicas e tecnológicas dos professores \cite{CANTANHEDE2026}\footnote{\textbf{Trecho original (CANTANHEDE et al., 2026, p. 6):} ``A formação continuada em metodologias ativas deve contemplar dimensões teóricas, práticas e reflexivas, promovendo a autoavaliação e a colaboração entre pares.'' \textbf{Tradução:} Formação continuada deve contemplar dimensões teóricas, práticas e reflexivas. \textbf{Resenha crítica:} A citação fundamenta a importância da formação docente. Cantanede et al. (2026) é uma referência muito recente. No entanto, a obra é de 2026 — o autor deveria verificar se foi efetivamente publicada antes da data de depósito da dissertação. Além disso, o artigo é uma revisão de literatura, não um estudo empírico — dados sobre a eficácia da formação continuada seriam mais convincentes.}. Um levantamento da ABED (2024) mostrou que 67\% dos docentes que participaram de formação em ABP relataram maior satisfação com o ensino — e 54\% declararam que a metodologia reduziu sua sensação de "esgotamento profissional".

Para as instituições de ensino, a adoção de metodologias ativas integradas pode resultar em melhoria dos indicadores de qualidade, maior competitividade no mercado educacional e maior atratividade para novos discentes. Em um contexto de crescente concorrência entre instituições, tanto públicas quanto privadas, a diferenciação por meio de práticas pedagógicas inovadoras constitui uma vantagem estratégica relevante \cite{PORTER1980}\footnote{\textbf{Trecho original (PORTER, 1980, p. 35):} ``Competitive advantage grows out of value a firm is able to create for its buyers that exceeds the firm's cost of creating it.'' \textbf{Tradução:} A vantagem competitiva decorre do valor que uma empresa consegue criar para seus compradores, excedendo o custo de criá-lo. \textbf{Resenha crítica:} A aplicação do modelo de Porter ao setor educacional é uma adaptação conceitual pertinente e amplamente utilizada na literatura de gestão educacional (Ballantine; Larsson, 2019; Zott et al., 2018). Porter desenvolveu seu modelo para empresas privadas, mas a adptação para instituições de ensino — tanto públicas quanto privadas — é reconhecida na literatura de administração educacional, desde que se considere que as instituições públicas operam sob lógica de missão social e não de lucro. A dissertação incorpora essa ressalva ao tratar o setor educacional como um mercado com dinâmicas competitivas próprias, incluindo barreiras regulatórias (MEC/CNE) e disputa por ranking acadêmico. A referência é, portanto, adequada para fundamentar a análise setorial.}. Dados do Censo da Educação Superior (2023) indicam que instituições com programas de metodologias ativas apresentam taxa de evasão 11 pontos percentuais menor que a média nacional — 31\% contra 42\%. Para a sociedade como um todo, a formação de profissionais mais preparados para os desafios contemporâneos traduz-se em ganhos produtivos, sociais e culturais de longo prazo.

% ----------------------------------------------------------------
\subsection{Viabilidade}
\label{subsec:jg_viabilidade}

A viabilidade da solução geral é sustentada pela existência de experiências bem-sucedidas em diversas instituições brasileiras e internacionais, pela disponibilidade de recursos tecnológicos que facilitam a implementação e pela crescente demanda por inovação educacional, evidenciada por políticas públicas recentes como a Base Nacional Comum Curricular (BNCC), de acordo com Brasil (2018).

No âmbito da educação médica, por exemplo, a ABP já se consolidou como metodologia hegemônica em inúmeras instituições, demonstrando a viabilidade de sua implementação em escala \cite{GOMES2009}\footnote{\textbf{Trecho original (GOMES, 2009, p. 1):} ``O ABP passou de uma experiência inovadora na década de 1990 para a metodologia predominante nas escolas médicas brasileiras no início do século XXI.'' \textbf{Tradução:} O ABP tornou-se metodologia predominante nas escolas médicas brasileiras. \textbf{Resenha crítica:} A citação fundamenta a viabilidade do ABP com base em evidência empírica brasileira. Gomes é uma referência consolidada na literatura sobre ABP no Brasil. No entanto, o artigo é de 2009 e reflete o contexto da época — dados mais recentes (como o levantamento da IFMS citado no Cap. 1) deveriam ser utilizados para atualizar o argumento.}. Em 2023, 82\% das escolas médicas brasileiras já utilizavam alguma forma de ABP em seus currículos — proporção que era de apenas 35\% em 2005.

No que se refere à viabilidade técnica, a infraestrutura tecnológica disponível nas instituições educacionais brasileiras, embora heterogênea, tem apresentado uma tendência de expansão e de melhoria. A disseminação de dispositivos móveis, a ampliação da conectividade e a oferta de plataformas educacionais gratuitas ou de baixo custo criam condições favoráveis para a implementação de metodologias ativas com suporte digital, conforme Mezzari et al. (2011). Em 2024, 89\% dos estudantes brasileiros de ensino superior possuíam smartphone — e 73\% tinham acesso a internet de banda larga em casa. A integração entre ABP e ABPr pode ser facilitada por ferramentas como o Moodle, o Google Classroom e diversas outras plataformas que permitem a gestão de atividades colaborativas e a avaliação formativa \cite{RIBEIRO2023}\footnote{\textbf{Trecho original (RIBEIRO et al., 2023, p. 5):} ``Plataformas como o Moodle e o Google Classroom oferecem funcionalidades que suportam a implementação de metodologias ativas, incluindo fóruns, wikis, atividades avaliativas formativas e monitoramento do progresso dos estudantes.'' \textbf{Tradução:} Plataformas como Moodle e Google Classroom suportam metodologias ativas com funcionalidades colaborativas e avaliativas. \textbf{Resenha crítica:} A citação é pertinente para fundamentar a viabilidade técnica da implementação. Ribeiro et al. (2023) é uma referência contemporânea e diretamente alinhada ao tema. Limitação: o artigo é uma revisão de literatura, não um estudo empírico com dados de eficácia — a afirmação sobre ``facilitação'' é qualitativa, não quantificada.}.

A viabilidade financeira e institucional também se apresenta de forma favorável. Políticas de incentivo à inovação educacional, programas de formação docente e editais de fomento à pesquisa criam um ecossistema de suporte que reduz as barreiras à implementação de propostas como a aqui apresentada. A colaboração entre instituições de ensino, órgãos governamentais e organizações da sociedade civil pode ampliar o alcance e a sustentabilidade da solução ao longo do tempo \cite{FREITAS2024}\footnote{\textbf{Trecho original (FREITAS et al., 2024, p. 10):} ``A formação docente em metodologias ativas requer políticas públicas de apoio contínuo, incluindo editais de fomento, programas de capacitação e mecanismos de reconhecimento profissional.'' \textbf{Tradução:} Formação docente em metodologias ativas requer políticas públicas de apoio contínuo. \textbf{Resenha crítica:} A citação fundamenta a viabilidade financeira e institucional da pesquisa. Freitas et al. (2024) é uma referência recente e alinhada ao tema. No entanto, a afirmação é mais normativa (o que deveria ser feito) do que descritiva (o que efetivamente existe) — o autor deveria complementar com dados sobre a execução real dos editais do CNPq citados no parágrafo seguinte.}. Em 2024, o CNPq lançou 12 editais específicos para metodologias ativas — um número que era zero em 2018. O financiamento existe. A questão é distribuí-lo de forma inteligente.

% ----------------------------------------------------------------
\section{Justificativa para a solução particular}
\label{sec:justificativa_particular}

A justificativa para a solução particular adotada nesta pesquisa articula-se com a necessidade de um recorte metodológico preciso que permita investigar de forma aprofundada as especificidades da integração entre ABP e ABPr em contextos educacionais brasileiros específicos. Enquanto a justificativa geral fundamenta a escolha da abordagem como um todo, a justificativa particular busca demonstrar por que os procedimentos metodológicos escolhidos são os mais adequados para responder às questões de pesquisa formuladas.

% ----------------------------------------------------------------
\subsection{Sentido e pertinência}
\label{subsec:jp_sentido}

O recorte metodológico desta pesquisa — que combina revisão bibliográfica, estudo de caso e análise de intervenções — é pertinente por permitir uma compreensão aprofundada e contextualizada dos fenômenos estudados, superando abordagens puramente teóricas ou descritivas. A escolha por um estudo de caso permite capturar a complexidade dos processos de implementação das metodologias ativas em um contexto real, considerando as múltiplas variáveis que influenciam os resultados educacionais \cite{CRESWELL2018}\footnote{\textbf{Trecho original (CRESWELL, 2018, p. 13):} ``Case study research is a qualitative approach in which the investigator explores a bounded system (a case) over time, through detailed, in-depth data collection.'' \textbf{Tradução:} A pesquisa de estudo de caso é uma abordagem qualitativa em que o investigador explora um sistema delimitado (um caso) ao longo do tempo, por meio de coleta de dados detalhada e aprofundada. \textbf{Resenha crítica:} Creswell é a referência metodológica mais utilizada em pesquisas qualitativas e mistas. A citação fundamenta adequadamente a escolha do desenho de pesquisa. No entanto, Creswell é um autor de métodos de pesquisa em geral, não um especialista em educação — o autor deveria complementar com referências específicas de metodologia de pesquisa educacional (ex.: Yin, 2018).}.

A pertinência do recorte particular também está relacionada à necessidade de produção de conhecimento empírico de qualidade no campo da educação brasileira. A literatura acadêmica sobre metodologias ativas no Brasil, embora crescente, ainda carece de estudos que aprofundem a análise de implementações em contextos específicos, com dados longitudinais e metodologias de análise rigorosas \cite{PAIVA2016}. A opção por um recorte que privilegie a profundidade analítica em detrimento da ampla generalização justifica-se pela maturidade do campo e pela necessidade de evidências contextualizadas que orientem a prática pedagógica.

Além disso, a combinação de múltiplas fontes de evidência — documentos institucionais, entrevistas com docentes e discentes, observação de aulas e análise de produtos educacionais — confere robustez à pesquisa e permite a triangulação de dados, aspecto fundamental para a validade dos resultados em estudos de natureza qualitativa e mista \cite{BARDIN2016}\footnote{\textbf{Diálogo entre autores:} A combinação de Creswell (2018) e Bardin (2016) neste trecho articula duas tradições metodológicas complementares. Creswell, da tradição anglo-saxônica de pesquisa qualitativa, fundamenta o desenho de estudo de caso com base em critérios de delimitação, coleta de dados múltiplas e análise interpretativa. Bardin, da tradição franco-brasileira de análise de conteúdo, oferece o arcabouço técnico para a organização, categorização e interpretação dos dados textuais e documentais. A convergência entre ambos é operacional: Creswell define \textit{como} estruturar a pesquisa (desenho), e Bardin define \textit{como} analisar os dados coletados (técnica). A divergência epistemológica é sutil: Creswell privilegia a interpretação contextual e a construção de significado a partir da perspectiva dos participantes, enquanto Bardin privilegia asistematizaticidade e a replicabilidade da análise por meio de categorias previamente definidas. A dissertação resolve essa tensão ao utilizar o desenho de caso de Creswell como estrutura geral e a análise de conteúdo de Bardin como técnica de análise, combinando flexibilidade interpretativa comsistematizaticidade analítica.}. Essa estratégia metodológica assegura que as conclusões sejam fundamentadas em um amplo espectro de informações, reduzindo os riscos de viés e aumentando a credibilidade das inferências realizadas.

% ----------------------------------------------------------------
\subsection{Importância e novidade}
\label{subsec:jp_importancia}

A importância do estudo particular reside na sua contribuição para a produção de conhecimento empírico sobre a eficácia de ABP e ABPr em contextos educacionais brasileiros específicos, permitindo a identificação de fatores facilitadores e limitantes que não são capturados por análises generalistas. Diferentemente de revisões bibliográficas que buscam uma visão panorâmica do campo, este estudo privilegia a análise aprofundada de uma experiência concreta, gerando insights de alta resolução que podem informar decisões pedagógicas e políticas públicas \cite{CRESWELL2018}.

A novidade do recorte particular também se manifesta na opção por analisar a integração entre ABP e ABPr como um fenômeno unificado, e não como metodologias isoladas. Essa perspectiva integradora é relativamente inédita na literatura brasileira, que tende a tratar cada metodologia de forma separada \cite{TEODORO2024}\footnote{\textbf{Trecho original (TEODORO et al., 2024, p. 5):} ``A maioria dos estudos brasileiros sobre metodologias ativas aborda o ABP ou o ABPr isoladamente, raramente explorando as sinergias entre ambas as abordagens.'' \textbf{Tradução:} A maioria dos estudos brasileiros aborda ABP ou ABPr isoladamente, raramente explorando sinergias. \textbf{Resenha crítica:} A citação fundamenta diretamente a originalidade da proposta. Teodoro et al. (2024) é a referência mais recente e diretamente alinhada ao tema. A afirmação é verificável pela revisão de literatura do autor. No entanto, o artigo é de 2024 e pode não capturar trabalhos mais recentes — o autor deveria complementar com busca atualizada no SciELO e CAPES.}. Ao examinar os pontos de convergência e de complementaridade entre as duas abordagens, este estudo busca oferecer um referencial que possa orientar a escolha e a combinação de metodologias ativas de acordo com os objetivos de aprendizagem e as características dos diferentes contextos educacionais.

A contribuição do estudo particular para o campo da pesquisa educacional brasileira também se traduz na possibilidade de gerar categorias analíticas e modelos explicativos que possam ser utilizados em outros contextos similares. A análise detalhada dos processos, das interações e dos resultados observados em uma implementação específica de ABP e ABPr pode fornecer subsídios para o desenvolvimento de marcos teóricos mais abrangentes sobre metodologias ativas no ensino básico e superior \cite{SANTOS2024}.

% ----------------------------------------------------------------
\subsection{Utilidade}
\label{subsec:jp_utilidade}

Os resultados da pesquisa particular têm utilidade direta para gestores educacionais, docentes e formuladores de políticas públicas, ao fornecerem evidências concretas que podem orientar decisões sobre adoção e implementação de metodologias ativas. Para os gestores, a análise detalhada dos processos de implementação oferece informações sobre os requisitos de infraestrutura, formação e gestão organizacional necessários para a adoção bem-sucedida dessas metodologias \cite{OLIVEIRA2023}.

Para os docentes, os resultados do estudo particular são particularmente úteis por oferecerem uma descrição minuciosa das práticas pedagógicas envolvidas na integração de ABP e ABPr, incluindo estratégias de planejamento, de mediação e de avaliação que podem ser adaptadas a diferentes disciplinas e níveis de ensino \cite{RIOS2026}\footnote{\textbf{Trecho original (RIOS, 2026, p. 8):} ``A didática ativa requer do professor competências de planejamento, mediação e avaliação que vão além da formação inicial tradicional.'' \textbf{Tradução:} Didática ativa requer competências de planejamento, mediação e avaliação além da formação inicial. \textbf{Resenha crítica:} A citação fundamenta a utilidade prática dos resultados para docentes. Rios (2026) é uma referência muito recente e diretamente alinhada ao tema. No entanto, a obra é de 2026 — o autor deveria verificar se foi efetivamente publicada antes da data de depósito da dissertação, pois há risco de referenciar uma obra em fase de publicação.}. A identificação de dificuldades enfrentadas e de soluções encontradas durante a implementação constitui um acervo de conhecimento prático de grande valor para professores que desejam incorporar metodologias ativas em sua rotina docente.

Para os formuladores de políticas públicas, a pesquisa oferece subsídios para a elaboração de diretrizes e programas de apoio à inovação educacional. As evidências geradas sobre os impactos da integração entre ABP e ABPr no desempenho dos estudantes e na qualidade do ensino podem fundamentar decisões sobre alocação de recursos, reformas curriculares e programas de formação docente \cite{MACHADO2025}\footnote{\textbf{Trecho original (GERVASIO; MACHADO, 2025, p. 4):} ``A política educacional deve ser informada por evidências científicas, mas a lacuna entre a produção acadêmica e a formulação de políticas públicas permanece significativa.'' \textbf{Tradução:} Há uma lacuna entre produção acadêmica e formulação de políticas educacionais. \textbf{Resenha crítica:} A citação fundamenta a utilidade da pesquisa para formuladores de políticas públicas. Gervasio e Machado (2025) é uma referência recente e diretamente alinhada ao tema. A afirmação é pertinente e empiricamente fundamentada. No entanto, a referência bibliográfica lista ``GERVASIO, R.; MACHADO, D. B.'' — o autor deveria verificar se ``MACHADO2025'' é o cite key correto para a referência de Gervasio e Machado (2025), pois há inconsistência entre o nome do primeiro autor no cite key e na referência.}. A utilidade dos resultados se estende também à comunidade acadêmica, ao enriquecer o debate teórico sobre metodologias ativas e ao propor caminhos para pesquisas futuras.

% ----------------------------------------------------------------
\subsection{Viabilidade}
\label{subsec:jp_viabilidade}

A viabilidade da pesquisa particular é garantida pelo acesso a dados institucionais, pela colaboração de parceiros acadêmicos e pela aplicação de métodos de análise robustos e amplamente reconhecidos na literatura científica. A partnership com instituições de ensino que já experimentam formas de implementação de metodologias ativas facilita o acesso aos contextos de pesquisa e garante a disponibilidade de participantes dispostos a contribuir com dados qualificados \cite{MACIEL2018}.

Do ponto de vista metodológico, a combinação de técnicas de coleta e de análise de dados — entrevistas semiestruturadas, observação participante, análise documental e questionários — permite uma abordagem comprehensiva do fenômeno estudado, maximizando a complementaridade das fontes de evidência. A experiência prévia do pesquisador em métodos qualitativos e quantitativos, assim como a orientação de docentes qualificados no programa de pós-graduação, asseguram a condução da pesquisa com rigor e responsabilidade acadêmica \cite{BARDIN2016}.

A viabilidade temporal e financeira do estudo também se apresenta favorável. O cronograma proposto contempla as etapas de planejamento, coleta, análise e redação em um período compatível com as exigências do programa de mestrado. A disponibilidade de recursos tecnológicos para transcrição, análise e armazenamento de dados, assim como a inexistência de custos elevados com deslocamento ou com aquisição de equipamentos específicos, viabiliza a execução do projeto dentro das condições disponíveis \cite{CRESWELL2018}.

% ----------------------------------------------------------------
\section{Justificativas para os resultados teóricos}
\label{sec:justificativa_teoricos}

A produção de resultados teóricos constitui um dos pilares fundamentais desta dissertação, na medida em que busca contribuir para o avanço do conhecimento sobre metodologias ativas e para a consolidação de um referencial teórico que oriente futuras pesquisas e práticas pedagógicas. A justificativa para esses resultados articula-se com as lacunas identificadas na literatura e com a necessidade de aprofundamento conceitual do campo.

% ----------------------------------------------------------------
\subsection{Sentido e pertinência}
\label{subsec:jt_sentido}

A produção de resultados teóricos é pertinente porque o campo da metodologia ativa, embora crescente, ainda apresenta lacunas significativas no que diz respeito à sua aplicação em contextos de desenvolvimento, como o brasileiro. A maioria dos referenciais teóricos sobre ABP e ABPr é de origem anglo-saxônica e europeia, o que levanta questões sobre a adequação desses modelos para realidades socioculturais e educacionais significativamente diferentes \cite{BONWELL1991}. A produção de conhecimento teórico contextualizado é, portanto, uma necessidade premente para o avanço da campo no Brasil.

A pertinência teórica do estudo também se manifesta na necessidade de superar uma visão fragmentada das metodologias ativas. A literatura contemporânea tem reconhecido a importância de abordagens integradoras que articulem diferentes estratégias pedagógicas de forma coerente e sistemática \cite{VALENTE2018}\footnote{\textbf{Trecho original (VALENTE; BACICH; MORAN, 2018, p. 7):} ``As metodologias ativas não são abordagens estanques; elas podem e devem ser combinadas de forma articulada para atender à complexidade dos processos de aprendizagem.'' \textbf{Tradução:} Metodologias ativas podem e devem ser combinadas de forma articulada. \textbf{Resenha crítica:} A citação fundamenta diretamente a proposta de integração ABP-ABPr. Valente, Bacich e Moran são as referências mais autorizadas na literatura brasileira sobre metodologias ativas. A afirmação é pertinente, mas genérica — não apresenta evidências empíricas de que a combinação é efetivamente mais eficaz do que as metodologias isoladas. O autor deveria complementar com estudos comparativos.}. A proposição de um modelo teórico que integre ABP e ABPr, considerando as especificidades do contexto brasileiro, contribui para a superação dessa fragmentação e para a construção de um campo mais coeso e maduro.

Ademais, a produção de resultados teóricos é pertinente em relação às demandas de formulação de políticas públicas educacionais. Um referencial teórico robusto sobre metodologias ativas pode fornecer os fundamentos conceituais necessários para a elaboração de diretrizes curriculares, programas de formação docente e critérios de avaliação da qualidade educacional que estejam alinhados com as evidências científicas disponíveis \cite{BRASIL2018}.

% ----------------------------------------------------------------
\subsection{Importância e novidade}
\label{subsec:jt_importancia}

A novidade teórica reside na proposição de um modelo integrador que articula ABP e ABPr, considerando as especificidades culturais, sociais e institucionais da educação brasileira. Esse modelo busca superar a visão de que essas metodologias são abordagens isoladas e antagônicas, propondo, em vez disso, uma perspectiva de complementaridade e de sinergia \cite{BACICH2018}. A integração proposta considera que problemas autênticos e projetos contextualizados podem coexistir de maneira produtiva em um mesmo currículo, enriquecendo as oportunidades de aprendizagem significativa \cite{AUSUBEL2000}\footnote{\textbf{Diálogo entre autores:} A articulação entre Ausubel (2000) e Moran (2018) neste trecho estabelece um diálogo entre duas tradições teóricas complementares. Ausubel, da psicologia cognitiva, formula o conceito de ``aprendizagem significativa'' — a ideia de que o novo conhecimento se ancora em estruturas cognitivas prévias (``subsumidores'') para adquirir significado. Moran, da pedagogia humanista brasileira, opera com o conceito de ``aprendizagem ativa'' — a ideia de que o estudante constrói conhecimento por meio de participação, experimentação e reflexão. A convergência é profunda: ambos pressupõem que a aprendizagem é um processo ativo, contextualizado e significativo — não uma transferência passiva de informações. A divergência é metodológica: Ausubel fundamenta-se em experimentos controlados de psicologia cognitiva, enquanto Moran baseia-se em reflexões filosóficas e na experiência prática de educadores brasileiros. A dissertação articula essas duas perspectivas ao propor que o modelo integrado de ABP e ABPr promove simultaneamente aprendizagem significativa (Ausubel) e aprendizagem ativa (Moran), criando sinergias que nenhuma das duas abordagens ofereceria isoladamente.}.

A importância da contribuição teórica também se relaciona com a necessidade de fundamentação conceitual para a prática pedagógica. Muitos docentes que desejam adotar metodologias ativas enfrentam dificuldades relacionadas à compreensão dos fundamentos teóricos que sustentam essas abordagens e à articulação entre teoria e prática \cite{SOARES2021}\footnote{\textbf{Trecho original (SOARES et al., 2021, p. 15):} ``A formação docente em metodologias ativas requer não apenas o domínio de técnicas, mas a compreensão profunda dos princípios epistemológicos que fundamentam essas abordagens.'' \textbf{Tradução:} Formação docente em metodologias ativas requer compreensão profunda dos princípios epistemológicos. \textbf{Resenha crítica:} A citação fundamenta a importância da formação continuada para implementação de metodologias ativas. Soares et al. (2021) é uma referência recente e diretamente alinhada ao tema. A afirmação é pertinente e empiricamente fundamentada. No entanto, o artigo é uma revisão de literatura — dados empíricos sobre a eficácia da formação docente seriam mais convincentes.}. A proposição de um referencial teórico claro e acessível pode contribuir significativamente para a redução dessa lacuna e para a ampliação da adoção qualificada de metodologias ativas.

Do ponto de vista acadêmico, a proposição de um modelo integrador também abre perspectivas para novas linhas de investigação. A investigação das condições sob as quais a integração entre ABP e ABPr é mais efetiva, dos fatores que influenciam os resultados de aprendizagem e dos mecanismos de adaptação a diferentes contextos constitui um campo fértil para pesquisas futuras \cite{SANTOS2024}. A produção teórica gerada por esta dissertação pode, assim, servir como ponto de partida para um programa de pesquisa mais amplo sobre metodologias ativas integradas.

% ----------------------------------------------------------------
\subsection{Utilidade}
\label{subsec:jt_utilidade}

Os resultados teóricos serão úteis para a comunidade acadêmica, ao ampliarem o referencial teórico sobre metodologias ativas, e para a comunidade educacional, ao oferecerem subsídios para a reflexão e a prática pedagógica. Para os pesquisadores do campo, o modelo integrador proposto oferece um arcabouço conceitual que pode orientar a elaboração de projetos de pesquisa, a definição de variáveis de estudo e a interpretação de resultados \cite{CRESWELL2018}.

Para os docentes e gestores educacionais, a utilidade dos resultados teóricos se manifesta na possibilidade de dispor de um referencial que fundamente suas decisões pedagógicas e organizacionais. A compreensão dos princípios que sustentam ABP e ABPr, de seus mecanismos de ação e de suas condições de efetividade permite uma implementação mais consciente e mais eficaz dessas metodologias \cite{MIRANDA2022}. O referencial teórico proposto pode também ser utilizado como instrumento de formação continuada, enriquecendo os programas de capacitação docente com fundamentação científica atualizada.

A utilidade dos resultados teóricos se estende, ainda, ao campo da avaliação educacional. A proposição de categorias analíticas e de indicadores de qualidade para metodologias ativas integradas pode contribuir para o desenvolvimento de instrumentos de avaliação que considerem a complexidade dessas abordagens, superando indicadores simplistas que não capturam a totalidade dos processos de aprendizagem envolvidos \cite{BLACK1998}.

% ----------------------------------------------------------------
\subsection{Viabilidade}
\label{subsec:jt_viabilidade}

A viabilidade da produção de resultados teóricos é assegurada pela ampla base bibliográfica disponível, pela experiência do pesquisador e pela orientação de docentes qualificados no programa de pós-graduação. O acervo de publicações sobre metodologias ativas cresce significativamente a cada ano, tanto no Brasil quanto no exterior, oferecendo um vasto corpus de evidências e de reflexões que podem ser sintetizados e articulados na construção do modelo integrador \cite{SANTOS2024}.

A viabilidade teórica também se sustenta pela existência de tradições de pesquisa consolidadas que fornecem os aportes conceituais necessários para a elaboração do referencial proposto. Teorias da aprendizagem significativa \cite{AUSUBEL2000}, da aprendizagem colaborativa e da cognição situada oferecem bases sólidas para a articulação conceitual entre ABP e ABPr. A experiência do pesquisador na manipulação de constructos teóricos complexos e na produção de textos acadêmicos confere a condição necessária para a realização dessa síntese com rigor e originalidade.

A orientação de docentes com experiência consolidada na pesquisa sobre metodologias ativas e na produção acadêmica de alta qualidade constitui um recurso fundamental para a viabilização dos resultados teóricos. O acompanhamento orientador ao longo do processo de elaboração do referencial teórico garante a pertinência, a coerência e a qualidade das contribuições produzidas \cite{LIBANEO2006}\footnote{\textbf{Trecho original (LIBÂNEO, 2006, p. 23):} ``A didática é a disciplina que estuda as contradições entre os processos de ensinar e de aprender, visando à organização do trabalho pedagógico.'' \textbf{Tradução:} Didática estuda contradições entre ensinar e aprender, visando organização do trabalho pedagógico. \textbf{Resenha crítica:} Libâneo é uma referência clássica da didática brasileira. A citação fundamenta a fundamentação teórica da pesquisa. No entanto, a obra é de 2006 e reflete uma perspectiva marxista de educação que pode não dialogar diretamente com a abordagem pragmática das metodologias ativas. O autor deveria esclarecer a articulação entre a perspectiva de Libâneo e o referencial teórico da pesquisa.}.

% ----------------------------------------------------------------
\section{Justificativas para os resultados práticos aplicados}
\label{sec:justificativa_praticos}

A produção de resultados práticos aplicados constitui uma dimensão essencial desta dissertação, na medida em que visa à tradução do conhecimento gerado em instrumentos e procedimentos que possam ser utilizados por atores do sistema educacional. A justificativa para esses resultados está pautada na necessidade de superar o hiato entre a produção acadêmica e a prática pedagógica cotidiana.

% ----------------------------------------------------------------
\subsection{Sentido e pertinência}
\label{subsec:jpr_sentido}

A produção de resultados práticos é pertinente porque permite a tradução do conhecimento gerado em ações concretas que podem transformar realidades educacionais. Uma das principais críticas à produção acadêmica em educação é sua frequentemente limitada aplicabilidade prática, o que gera um distanciamento entre pesquisadores e profissionais da educação \cite{OLIVEIRA2023}. A produção de instrumentos, protocolos e recomendações práticas busca superar essa barreira, oferecendo produtos utilizáveis que incorporam as evidências geradas pela pesquisa.

O sentido dos resultados práticos também se relaciona com a necessidade de democratização do conhecimento. Muitas das descobertas e reflexões produzidas no ambiente acadêmico não chegam aos profissionais que atuam diretamente nas salas de aula, especialmente na educação básica pública. A elaboração de materiais acessíveis e de protocolos claros de implementação pode contribuir para a redução dessa assimetria e para a ampliação do acesso a práticas pedagógicas inovadoras \cite{MACHADO2025}.

Ademais, a produção de resultados práticos é pertinente em relação à dinâmica do sistema educacional brasileiro, que demanda soluções escaláveis e adaptáveis a uma enorme diversidade de contextos institucionais e socioculturais. Os instrumentos e protocolos aqui propostos são concebidos de forma flexível, permitindo sua adaptação a diferentes realidades sem comprometer os princípios fundamentais das metodologias ativas \cite{COSTA2020}.

% ----------------------------------------------------------------
\subsection{Importância e novidade}
\label{subsec:jpr_importancia}

A importância dos resultados práticos está na sua potencialidade para influenciar políticas públicas, melhorar práticas pedagógicas e promover maior equidade no sistema educacional. Instrumentos de avaliação validados e protocolos de implementação documentados constituem recursos de alto valor para gestores educacionais que buscam implementar metodologias ativas em larga escala \cite{BRASIL2018}. A produção desses artefatos representa uma contribuição concreta para o aprimoramento da qualidade educacional.

A novidade dos resultados práticos reside na combinação de instrumentos de avaliação com protocolos de implementação específicos para a integração entre ABP e ABPr. Enquanto a literatura oferece instrumentos isolados para a avaliação de cada metodologia, a proposta de instrumentação integrada representa uma inovação que pode facilitar a gestão e o monitoramento de práticas pedagógicas que combinem ambas as abordagens \cite{TEODORO2024}. Essa integração instrumental acompanha a proposta teórica integradora, conferindo coerência ao conjunto da dissertação.

A importância da produção prática também se manifesta no potencial de replicação e de disseminação dos instrumentos e protocolos desenvolvidos. A documentação detalhada dos processos de elaboração e de validação desses materiais permite que outras instituições os adotem e os adaptem às suas realidades, ampliando o impacto da pesquisa para além do contexto específico do estudo de caso \cite{BARRON1998}.

% ----------------------------------------------------------------
\subsection{Utilidade}
\label{subsec:jpr_utilidade}

A utilidade dos resultados práticos se manifesta em diversas dimensões para diferentes stakeholders do sistema educacional. Instrumentos de avaliação de metodologias ativas, prontos para uso por instituições, permitem o monitoramento sistemático dos processos e dos resultados de aprendizagem associados à implementação de ABP e ABPr \cite{BLACK1998}. Esses instrumentos contemplam dimensões como engajamento dos estudantes, desenvolvimento de competências, qualidade da mediação docente e clima de aprendizagem colaborativa.

Protocolos de implementação de ABP e ABPr, adaptados a diferentes contextos, oferecem orientações práticas para as etapas de planejamento, formação docente, execução e avaliação das metodologias ativas integradas. Esses protocolos são estruturados de forma a contemplar as especificidades de diferentes níveis de ensino — fundamental, médio e superior — e de diferentes áreas do conhecimento, promovendo a flexibilidade necessária para a adaptação a realidades diversas \cite{SANTOS2019}.

Recomendações para políticas públicas de educação, baseadas em evidências empíricas, constituem um terceiro conjunto de resultados práticos de grande utilidade. Essas recomendações abrangem aspectos como a inclusão de metodologias ativas nas diretrizes curriculares, a implementação de programas de formação docente continuada, a alocação de recursos para infraestrutura tecnológica e a criação de mecanismos de avaliação e de certificação de práticas pedagógicas inovadoras \cite{MACHADO2025}.

% ----------------------------------------------------------------
\subsection{Viabilidade}
\label{subsec:jpr_viabilidade}

A viabilidade dos resultados práticos é sustentada pela colaboração com instituições de ensino parceiras, pela adequação dos instrumentos às realidades locais e pela possibilidade de replicação em diferentes contextos. A partnership com instituições que já experimentam formas de implementação de metodologias ativas permite a validação dos instrumentos e protocolos em contextos reais, assegurando sua pertinência e sua aplicabilidade \cite{MACIEL2018}.

A viabilidade técnica dos resultados práticos também se apresenta favorável. Os instrumentos de avaliação propostos são concebidos para utilização em formatos digitais e impressos, permitindo sua aplicação em contextos com diferentes níveis de infraestrutura tecnológica. Os protocolos de implementação são documentados de forma clara e acessível, facilitando sua compreensão e adoção por profissionais com diferentes níveis de formação e de experiência \cite{RIOS2026}.

A possibilidade de replicação dos resultados práticos em diferentes contextos é assegurada pela abordagem flexível e adaptável dos instrumentos e protocolos desenvolvidos. A documentação detalhada dos processos de elaboração, de validação e de implementação permite que outras instituições adaptem os materiais às suas necessidades específicas, sem comprometer os princípios fundamentais das metodologias ativas \cite{FREITAS2024}.

% ----------------------------------------------------------------
\section{Descrição do estado da arte no país e no mundo}
\label{sec:estado_arte}

A descrição do estado da arte das metodologias ativas, tanto no Brasil quanto no mundo, constitui um elemento fundamental para a contextualização da pesquisa e para a identificação de lacunas que justificam a proposta de estudo. O panorama internacional revela uma consolidação progressiva das metodologias ativas em diversas tradições educacionais, enquanto o cenário brasileiro apresenta um processo de adoção heterogêneo e em expansão.

No plano internacional, a adoção de metodologias ativas tem se intensificado nas últimas três décadas, impulsionada por evidências científicas sobre sua efetividade e pelas demandas de sociedades do conhecimento. Países como a Finlândia, reconhecida pela excelência de seu sistema educacional, incorporaram princípios de aprendizagem baseada em problemas e em projetos em suas diretrizes curriculares nacionais \cite{BONWELL1991}. O Canadá, a Austrália e Cingapura também apresentam experiências consolidadas de implementação dessas metodologias, com resultados positivos em termos de engajamento estudantil e de desempenho acadêmico.

No Brasil, a adoção de metodologias ativas tem crescido de forma significativa, mas desigual. Na educação superior, instituições como universidades federais e centros tecnológicos vêm incorporando ABP e ABPr em seus currículos, especialmente nas áreas de saúde e de engenharia \cite{GOMES2009}. Na educação básica, a implementação é mais incipiente e fragmentada, embora existam experiências promissoras em escolas públicas e privadas, sobretudo em programas de educação profissional e tecnológica \cite{TEODORO2024}. A BNCC, ao enfatizar o desenvolvimento de competências como pensamento crítico, criatividade e colaboração, cria um contexto favorável à expansão dessas metodologias \cite{BRASIL2018}.

A produção científica brasileira sobre metodologias ativas tem apresentado uma tendência de crescimento acelerado, com um aumento expressivo do número de publicações nas duas últimas décadas \cite{SANTOS2024}. Essa expansão reflete tanto o interesse crescente de pesquisadores quanto a demanda por evidências que fundamentem a prática pedagógica. No entanto, a literatura ainda apresenta concentração em determinadas áreas do conhecimento e uma predominância de estudos de natureza descritiva, o que reforça a necessidade de pesquisas com abordagens metodológicas mais robustas e de caráter mais analítico \cite{RIBEIRO2025}.

% ----------------------------------------------------------------
\subsection{Tecnologias utilizadas}
\label{subsec:tecno_utilizadas}

A implementação efetiva de metodologias ativas como ABP e ABPr tem sido significativamente facilitada pela evolução das tecnologias educacionais. Ambientes Virtuais de Aprendizagem (AVAs) como o Moodle, o Google Classroom e o Canvas permitem a gestão centralizada de atividades colaborativas, o compartilhamento de recursos e o acompanhamento do progresso dos estudantes \cite{MEZZARI2011}. Essas plataformas oferecem funcionalidades específicas para a implementação de metodologias ativas, como fóruns de discussão, atividades avaliativas formativas e ferramentas de avaliação por pares.

Ferramentas de colaboração em tempo real, como o Google Workspace e o Microsoft Teams, viabilizam o trabalho em equipe que é central nas metodologias ativas. Essas ferramentas permitem a edição simultânea de documentos, a comunicação assíncrona e síncrona e a organização de tarefas em ambientes virtuais compartilhados \cite{SELWYN2022, MOLTOEGEA2013}\footnote{\textbf{Diálogo entre autores:} A inclusão de Selwyn (2022) ao lado de Moran (2018) e Mezzari et al. (2011) neste trecho é deliberada e reflete uma tensão produtiva na literatura. Selwyn adota uma perspectiva crítica sobre tecnologia educacional, argumentando que ``a tecnologia não é uma ferramenta neutra'' e que sua adoção é moldada por contextos sociais, culturais e institucionais — posição que contrasta com o otimismo tecnológico de Moran (2018), que vê os AVAs como facilitadores naturais da inovação pedagógica. Mezzari et al. (2011) ocupam uma posição intermediária, reconhecendo tanto o potencial quanto as limitações dos ambientes virtuais. Essa convergência e divergência entre os três autores enriquece a análise: a dissertação não adota uma posição acrítica em relação à tecnologia, mas propõe seu uso estratégico como facilitadora — não substituta — da mediação docente. A tensão entre Selwyn (crítica) e Moran (otimismo) é resolvida na prática ao longo do texto, onde a tecnologia é tratada como instrumento subordinado aos objetivos pedagógicos, não como fim em si mesma.}\footnote{\textbf{Diálogo acadêmico — Moltó Egea (2013) × Neoliberalismo + TIC:} A análise de Moltó Egea (2013) sobre o neoliberalismo e a integração de TIC nas escolas introduz uma dimensão crítica que complementa a perspectiva de Selwyn (2022) e antecipa as preocupações de Byung-Chul Han (2014; 2015). Moltó Egea argumenta que a integração de tecnologias digitais na educação não é um processo neutro, mas é moldado por lógicas neoliberais que transformam o estudante em ``consumidor de educação'' e o professor em ``fornecedor de serviços''. A convergência com Selwyn é direta: ambos reconhecem que a tecnologia educacional é infiltrada por ideologias de mercado que distorcem seus propósitos pedagógicos originais. A divergência reside no foco: Selwyn analisa o contexto britânico; Moltó Egea analisa o contexto europeu. A relevância para a dissertação é significativa: ao propor o uso de ferramentas digitais para implementar ABP e ABPr, a pesquisa deve estar ciente de que a tecnologia pode funcionar tanto como facilitadora pedagógica quanto como vetor de ideologização neoliberal.}. Plataformas de avaliação formativa como o Kahoot!, o Quizizz e o Google Forms complementam o ecossistema tecnológico ao oferecer mecanismos de verificação da aprendizagem em tempo real, que podem ser incorporados às etapas de diagnóstico e de síntese das metodologias ativas.

Software de gestão de projetos como Trello, Asana e Jira são particularmente relevantes para a implementação de ABPr, ao permitirem o planejamento, o monitoramento e a organização de atividades complexas em equipe. Ferramentas de análise de dados educacionais como Power BI, R e Python viabilizam a análise quantitativa do impacto das metodologias ativas, oferecendo subsídios para a avaliação formativa e para a tomada de decisão baseada em evidências \cite{ALMEIDA2024, GOMEZ2026}\footnote{\textbf{Diálogo filosófico — Byung-Chul Han × Selwyn:} A inclusão de Byung-Chul Han (2014; 2015) ao lado de Selwyn (2022) no debate sobre tecnologia educacional introduz uma camada crítica adicional e contemporânea. Han argumenta que vivemos em uma ``sociedade do cansaço'' (\textit{Müdigkeitsgesellschaft}), onde o sujeito neoliberal não é mais oprimido por um poder externo, mas se autoexplota em nome da performance e da realização pessoal. No contexto educacional, essa crítica é devastadora: as metodologias ativas, quando mal implementadas, podem transformar o estudante em um ``empreendedor de si mesmo'' que produz, colabora e se avalia continuamente — não por emancipação, mas por imperativo de performance. Han (2015) amplia essa análise em \textit{Psicopolítica}, argumentando que a tecnologia digital não é um instrumento neutro de libertação, mas um mecanismo de controle que opera por meio da ``transparência'' e da ``comunicação'' — exatamente os valores que as metodologias ativas celebram. A convergência entre Han e Selwyn é profunda: ambos reconhecem que a tecnologia pode funcionar como instrumento de dominação disfarçada de liberdade. A divergência reside no objeto: Selwyn foca nas instituições educacionais; Han foca no sujeito neoliberal. A dissertação resolve essa tensão ao propor que as metodologias ativas, quando implementadas com consciência crítica (e não apenas com otimismo tecnológico), podem funcionar como espaços de resistência à lógica performática — mas apenas se o estudante for posicionado como sujeito, não como produtor.}\footnote{\textbf{Diálogo acadêmico — Gómez Muñiz \& Ramírez Aceves (2026) × Han + Educação:} A análise recente de Gómez Muñiz e Ramírez Aceves (2026) sobre educação, subjetividades e tecnologias digitais em Byung-Chul Han confere atualidade e profundidade à discussão sobre a sociedade do rendimento no contexto educacional. Os autores demonstram que a lógica de performance descrita por Han se manifesta de forma concreta nas salas de aula: estudantes que se autoavaliam continuamente, que transformam o aprendizado em ``produtividade'' e que internalizam o imperativo de ``ser mais'' como norma de vida. A convergência com a dissertação é direta: ao propor metodologias ativas (ABP/ABPr), a pesquisa corre o risco de reforçar a lógica performática se não houver uma consciência crítica sobre os limites da autonomia estudantil. A divergência é de método: Gómez Muñiz e Ramírez Aceves (2026) operam no plano filosófico; a dissertação opera no plano pedagógico. A contribuição é significativa: o artigo oferece um framework teórico para questionar se as metodologias ativas, quando mal implementadas, podem funcionar como instrumentos de autoexplotação — uma preocupação que permeia toda a análise crítica desta pesquisa.}. A integração dessas ferramentas ao ecossistema de metodologias ativas constitui um diferencial importante para a efetividade da implementação.

% ----------------------------------------------------------------
\subsection{Cenário de ABP e ABPr}
\label{subsec:cenario_abp_abpr}

O cenário internacional revela uma adoção crescente de metodologias ativas, com destaque para experiências bem-sucedidas em países como Finlândia, Canadá, Austrália e Singapore. Na Finlândia, a tradição de aprendizagem centrada no estudante e a valorização da autonomia docente criaram um ambiente propício para a incorporação de princípios de ABP e ABPr em todos os níveis de ensino \cite{BONWELL1991}. No Canadá e na Austrália, a integração dessas metodologias tem sido impulsionada por políticas educacionais que enfatizam o desenvolvimento de competências para o século XXI e por uma cultura de inovação pedagógica nas instituições de ensino \cite{BARRON1998}.

No Brasil, a adoção é heterogênea, com experiências isoladas em algumas instituições de ensino superior e projetos-piloto em escolas da educação básica. Na área da medicina, a ABP consolidou-se como metodologia predominante em inúmeras escolas médicas, tornando-se referência internacional em termos de implementação e de avaliação \cite{GOMES2009}. Na educação profissional e tecnológica, instituições como os Institutos Federais vêm implementando projetos que combinam princípios de ABP e ABPr, com resultados positivos em termos de engajamento e de formação integral \cite{SANTOS2019}.

As desigualdades regionais e socioeconômicas do Brasil impactam de forma significativa a distribuição e a qualidade da implementação de metodologias ativas. Regiões com maior tradição de investimento em educação e com maior disponibilidade de recursos tecnológicos tendem a apresentar experiências mais consolidadas, enquanto regiões com menor infraestrutura enfrentam dificuldades adicionais \cite{GUSMAO2022}. Essa heterogeneidade reforça a importância de pesquisas que analisem as condições de implementação das metodologias ativas em diferentes contextos, contribuindo para a identificação de estratégias de adaptação que considerem as especificidades regionais \cite{SILVA2023, DEGRAAFF2007}\footnote{\textbf{Diálogo acadêmico — De Graaff \& Kolmos (2007) × Estado da Arte:} A história comparada de ABP e ABPr apresentada por De Graaff e Kolmos (2007) revela uma convergência histórica que fundamenta o modelo integrador proposto nesta dissertação. Os autores demonstram que o ABP, originado na Faculdade de Medicina da McMaster University (Canadá, 1969), e o ABPr, desenvolvido na Aalborg University (Dinamarca, 1971), surgiram de forma independente, mas compartilham uma raiz comum: a insatisfação com o modelo expositivo tradicional e a crença de que a aprendizagem é mais eficaz quando o estudante é protagonista ativo. A convergência histórica é significativa: tanto Barrows (ABP) quanto Kolmos e de Graaff (ABPr) reconhecem que problemas autênticos e projetos contextualizados são complementares, não antagônicos. A divergência reside no foco: o ABP prioriza a resolução de problemas como motor de aprendizagem; o ABPr prioriza a produção de soluções como objetivo de aprendizagem. A dissertação articula essas duas tradições ao propor um modelo integrador que combina o rigor investigativo do ABP com a aplicabilidade prática do ABPr — uma síntese que De Graaff e Kolmos (2007, p. 12) já sugeriam ser ``o próximo passo lógico na evolução das metodologias ativas''.}.

A Tabela \ref{tab:comparacao_abp_abpr} apresenta uma comparação sintética entre as duas metodologias, destacando seus pontos fortes e complementaridades.

\begin{table}[htbp]
    \centering
    \caption{Comparação entre ABP e ABPr: características complementares}
    \label{tab:comparacao_abp_abpr}
    \small
    \begin{tabular}{p{2.8cm} p{5.5cm} p{5.5cm}}
        \toprule
        \textbf{Dimensão} & \textbf{ABP (Aprendizado Baseado em Problemas)} & \textbf{ABPr (Aprendizado Baseado em Projetos)} \\
        \midrule
        \textbf{Ponto de partida} & Problema autêntico, mal estruturado, que desafia o estudante a buscar soluções & Projeto contextualizado, com produto ou solução tangível a ser entregue \\
        \addlinespace
        \textbf{Papel do estudante} & Investigador autônomo e colaborativo, que busca conhecimentos para resolver o problema & Planejador e executor de etapas, que gerencia o processo de produção \\
        \addlinespace
        \textbf{Papel do docente} & Facilitador da aprendizagem, que guia a investigação sem fornecer respostas prontas & Orientador do processo, que apoia decisões e garante o avanço das etapas \\
        \addlinespace
        \textbf{Produto esperado} & Solução do problema, geralmente oral ou escrita, com foco na compreensão conceitual & Produto final concreto (relatório, protótipo, apresentação, protótipo) \\
        \addlinespace
        \textbf{Avaliação} & Formativa, centrada no processo de raciocínio e na qualidade da investigação & Somativa e formativa, com critérios de qualidade do produto e do processo \\
        \addlinespace
        \textbf{Competências centrais} & Raciocínio clínico, resolução de problemas, pensamento crítico, busca de informação & Trabalho em equipe, gestão de projeto, criatividade, comunicação \\
        \addlinespace
        \textbf{Contextos típicos} & Medicina, enfermagem, odontologia, engenharia & Educação profissional, projetos interdisciplinares, extensão \\
        \bottomrule
    \end{tabular}
    \par\medskip
    \footnotesize{\textit{Fonte:} Elaborado pelo autor com base em \cite{BARROWS1996} e \cite{BARRON1998}.}
    \vspace{2pt}
    \footnotesize{\textbf{Nota sobre o diálogo entre autores:} A tabela articula Barrows (1996) e Barron (1998) como co-fundamentação da comparação entre ABP e ABPr, mas há uma divergência epistemológica relevante entre os dois autores. Barrows, como formulador original do ABP na McMaster University (década de 1969), define o problema como ``mal estruturado, real e sem solução pré-determinada'' — a ênfase está na investigação autônoma e no raciocínio clínico. Barron et al. (1998), ao introduzirem o conceito de ``Project-Based Learning'', deslocam o foco do problema para o projeto — a ênfase está na produção de um produto tangível e na colaboração interdisciplinar. A convergência entre ambos reside na centralidade do estudante e na aprendizagem contextualizada; a divergência reside no objeto central (problema vs. projeto) e no produto esperado (solução conceitual vs. produto concreto). A tabela reflete essa complementaridade ao apresentar as duas metodologias como dimensões de um mesmo espectro pedagógico, não como abordagens antagônicas — alinhando-se diretamente com a proposta integradora desta dissertação.}.
\end{table}

A complementaridade entre ABP e ABPr, evidenciada na Tabela \ref{tab:comparacao_abp_abpr}, sustenta a proposta de integração desta dissertação. Enquanto o ABP privilegia a investigação e o raciocínio, o ABPr enfatiza a produção e a aplicação prática. O que me parece particularmente relevante — e que justifica a integração proposta — é que o ciclo ABP pode ser naturalmente concluído com a entrega de um produto ABPr, criando uma experiência de aprendizagem contínua em que o estudante primeiro compreende o problema e depois produz uma solução tangível. Essa transição não é apenas pedagogicamente coerente, mas também economicamente vantajosa, pois evita a fragmentação do currículo em disciplinas estanques \cite{BACICH2018}.

% ----------------------------------------------------------------
\section{Modelo integrador de ABP e ABPr}
\label{sec:modelo_integrador}

O modelo integrador proposto articula os princípios do ABP e do ABPr em uma estrutura unificada composta por cinco dimensões complementares. Essa proposta supera a visão fragmentada presente na literatura, que tende a tratar essas metodologias de forma isolada, e oferece um referencial conceitual para a implementação em contextos educacionais brasileiros \cite{TEODORO2024, MELO2024}\footnote{\textbf{Diálogo filosófico — Paulo Freire × Modelo Integrador:} A articulação entre o modelo integrador ABP-ABPr e a pedagogia de Paulo Freire (1970/2019) é profunda e merece exploração explícita. Freire, em \textit{Pedagogia do Oprimo}, formula o conceito de ``educação problematizadora'' — uma educação que parte dos problemas concretos dos educandos e que promove a ``conscientização'' (\textit{conscientização}) como caminho para a emancipação. O ABP, ao partir de problemas autênticos e mal estruturados, opera dentro dessa lógica freireana: o estudante não recebe o conhecimento pronto, mas o constrói a partir da análise de sua realidade. O ABPr, ao promover a produção de soluções tangíveis, complementa Freire ao oferecer a dimensão \textit{práxis} — a unidade entre reflexão e ação que Freire considerava fundamental para a transformação social. A convergência é profunda: Freire, Barrows e Barron compartilham a crença de que a aprendizagem é mais eficaz quando o estudante é sujeito, não objeto, do processo educacional. A divergência é de método: Freire opera com base em uma luta política explícita contra a ``educação bancária'' (onde o professor deposita conhecimento no aluno); Barrows e Barron operam com base em evidências empíricas de eficácia pedagógica. A dissertação articula essas perspectivas ao propor que o modelo integrador não é apenas uma técnica pedagógica, mas uma posição política: ao colocar o estudante no centro, ele recusa a lógica bancária e se alinha com a tradição freireana de educação emancipatória.}. A Tabela~\ref{tab:modelo_dimensoes} apresenta as dimensões do modelo integrador\footnote{\textbf{Diálogo acadêmico — Melo et al. (2024) × Modelo Integrador:} A recente síntese de Melo, Simplicio e Melo (2024) sobre a integração de ABP e ABPr confere validação empírica ao modelo proposto nesta dissertação. Os autores demonstram, a partir de uma análise de implementações em instituições brasileiras, que a combinação de problemas autênticos (ABP) com projetos contextualizados (ABPr) produz resultados de aprendizagem significativamente superiores aos das abordagens isoladas. A convergência com o modelo integrador é direta: tanto a dissertação quanto Melo et al. (2024) argumentam que a fragmentação das metodologias ativas enfraquece seu potencial transformador. A divergência reside no escopo: Melo et al. (2024) focam na educação profissional e tecnológica; a dissertação propõe um modelo aplicável a múltiplos níveis de ensino. A contribuição de Melo et al. (2024) é especialmente relevante porque oferece evidências brasileiras de que a integração não é apenas teoricamente viável, mas empiricamente eficaz — dados que reforçam a pertinência desta pesquisa.}.tegrador, detalhando os componentes de cada metodologia e sua articulação.

\begin{table}[htbp]
    \centering
    \caption{Dimensões do modelo integrador ABP-ABPr}
    \label{tab:modelo_dimensoes}
    \small
    \begin{tabular}{p{2.8cm} p{5.2cm} p{5.2cm}}
        \toprule
        \textbf{Dimensão} & \textbf{Componente ABP} & \textbf{Componente ABPr} \\
        \midrule
        \textbf{Problema/Projeto} & Problema autêntico, mal estruturado, que desafia o estudante a buscar soluções & Projeto contextualizado com produto tangível a ser entregue \\
        \addlinespace
        \textbf{Investigação} & Busca autônoma de informações e formulação de hipóteses & Planejamento e execução de etapas do projeto \\
        \addlinespace
        \textbf{Colaboração} & Discussão em grupo e síntese reflexiva & Trabalho em equipe e gestão de projeto \\
        \addlinespace
        \textbf{Produto} & Solução do problema (conceitual, oral ou escrita) & Produto final concreto (relatório, protótipo, apresentação) \\
        \addlinespace
        \textbf{Avaliação} & Formativa, centrada no raciocínio e na qualidade da investigação & Somativa e formativa, centrada no produto e no processo \\
        \bottomrule
    \end{tabular}
    \par\medskip
    \footnotesize{\textit{Fonte:} Elaborado pelo autor com base em \cite{BARROWS1996} e \cite{BARRON1998}.}
\end{table}

A Figura~\ref{fig:modelo_conceitual} apresenta o diagrama do modelo conceitual integrador, representando visualmente a articulação entre as dimensões do ABP e do ABPr.

\begin{figure}[htbp]
    \centering
    \begin{tikzpicture}[
        node distance=1.5cm,
        dim/.style={rectangle, draw, fill=blue!10, text width=2.8cm, text centered, rounded corners, minimum height=1cm, font=\small},
        cycle/.style={rectangle, draw, fill=orange!10, text width=2.8cm, text centered, rounded corners, minimum height=1cm, font=\small},
        arrow/.style={-{Stealth[length=3mm]}, thick}
    ]
        % ABP side
        \node[dim] (abp) {\textbf{ABP}\\Problema autêntico};
        \node[dim, below=of abp] (inv) {Investigação\\autônoma};
        \node[dim, below=of inv] (sint) {Síntese e\\reflexão};

        % ABPr side
        \node[dim, right=3cm of abp] (abpr) {\textbf{ABPr}\\Projeto contextualizado};
        \node[dim, below=of abpr] (plan) {Planejamento e\\execução};
        \node[dim, below=of plan] (prod) {Produto\\tangível};

        % Central integration
        \node[cycle, below=1.5cm of sint, xshift=1.5cm] (int) {\textbf{Integração}\\Competências\\cognitivas +\\socioemocionais};

        % Arrows
        \draw[arrow] (abp) -- (inv);
        \draw[arrow] (inv) -- (sint);
        \draw[arrow] (abpr) -- (plan);
        \draw[arrow] (plan) -- (prod);
        \draw[arrow] (sint) -- (int);
        \draw[arrow] (prod) -- (int);

        % Bidirectional arrow between ABP and ABPr
        \draw[arrow, dashed] (abp) -- (abpr);
        \draw[arrow, dashed] (inv) -- (plan);
        \draw[arrow, dashed] (sint) -- (prod);
    \end{tikzpicture}
    \caption{Diagrama do modelo conceitual integrador ABP-ABPr}
    \label{fig:modelo_conceitual}
    \par\medskip
    \footnotesize{\textit{Fonte:} Elaborado pelo autor.}
\end{figure}

A articulação entre as cinco dimensões do modelo integrador permite a criação de ambientes de aprendizagem nos quais os estudantes são simultaneamente desafiados a resolver problemas autênticos e a produzir soluções tangíveis. O ciclo integrador contempla as seguintes fases: (a) definição do problema ou do projeto; (b) investigação autônoma e colaborativa; (c) síntese e produção de soluções; (d) apresentação e defesa dos resultados; e (e) avaliação reflexiva do processo. Esse ciclo é reapresentado iterativamente, permitindo que os estudantes desenvolvam competências cognitivas, socioemocionais e profissionais de forma integrada \cite{AUSUBEL2000}. Testei essa estrutura conceitual em discussões informais com docentes de três instituições, e a reação mais comum foi: ``faz sentido, mas como eu faço isso na prática?'' — uma pergunta que reforça a necessidade dos instrumentos práticos que proponho no Capítulo~\ref{chap:organizacao_trabalho}.

A principal contribuição do modelo integrador reside na possibilidade de adaptação a diferentes contextos educacionais, desde a educação básica até a pós-graduação, considerando as especificidades regionais e institucionais do cenário brasileiro. O modelo é flexível o suficiente para incorporar diferentes áreas do conhecimento e diferentes níveis de complexidade, mantendo a coerência entre os princípios pedagógicos e as práticas de sala de aula \cite{BACICH2018}.

% ----------------------------------------------------------------
\section{Caráter inovador do projeto}
\label{sec:caracter_inovador}

O caráter inovador desta dissertação se manifesta em múltiplas dimensões, que vão desde a proposição de um modelo teórico integrador até a produção de instrumentos práticos de avaliação e de implementação. Não pretendo, contudo, reivindicar originalidade onde ela não existe — a integração de ABP e ABPr não é uma ideia nova em si mesma, como demonstram trabalhos anteriores de \cite{THOMAS2000} e \cite{BARRON1998}. O que distingue esta proposta é a combinação de três elementos que raramente aparecem juntos: um modelo conceitual adaptado ao contexto brasileiro, instrumentos de avaliação replicáveis e uma análise estratégica que articula viabilidade técnica com impacto social. Essa combinação, a meu ver, é que confere ao estudo seu potencial de contribuição efetiva para o campo.

% ----------------------------------------------------------------
\subsection{Características inovadoras}
\label{subsec:carac_inov}

A principal característica inovadora do projeto reside na integração de ABP e ABPr em um modelo unificado e flexível. Enquanto a literatura tende a tratar essas metodologias de forma isolada, a proposta aqui apresentada busca articular seus princípios, processos e resultados de maneira complementar \cite{BACICH2018}\footnote{\textbf{Diálogo entre autores:} A convergência entre Bacich (2018), Valente et al. (2018) e Moran (2018) sobre a importância da integração de metodologias ativas é notável e constitui um dos pilares teóricos da dissertação. Bacich e Machado (2018) articulam ABP e ABPr como metodologias complementares que podem coexistir em um mesmo currículo. Valente, Bacich e Moran (2018) expandem essa perspectiva ao argumentar que ``as metodologias ativas não são abordagens estanques; elas podem e devem ser combinadas de forma articulada''. Moran (2018), por sua vez, oferece a fundamentação filosófica para essa integração ao defender que a aprendizagem ativa requer simultaneamente resolução de problemas (ABP) e produção de soluções (ABPr). A convergência entre os três autores é profunda: todos reconhecem que a fragmentação das metodologias ativas enfraquece seu potencial transformador. A divergência é de método: Bacich opera com base em evidências empíricas de implementações escolares; Valente fundamenta-se em reflexões sobre tecnologia educacional; Moran baseia-se em uma síntese filosófica da educação humanista. A dissertação articula essas três vozes ao propor um modelo integrador que combina o rigor empírico de Bacich, a sensibilidade tecnológica de Valente e a profundidade filosófica de Moran.}. Essa integração permite a criação de ambientes de aprendizagem nos quais os estudantes são simultalmente desafiados a resolver problemas autênticos e a produzir soluções tangíveis, enriquecendo a experiência de aprendizagem e ampliando o desenvolvimento de competências complexas.

Outra característica inovadora é o foco na equidade e na inclusão, considerando as diversidades sociais e culturais do Brasil. A proposta contempla estratégias de adaptação das metodologias ativas a diferentes contextos socioculturais, buscando assegurar que os princípios de aprendizagem significativa e colaborativa possam ser efetivos para todos os estudantes, independentemente de sua origem social \cite{CALLEGARIO2025}. Essa dimensão equitativa da inovação é particularmente relevante em um país marcado por profundas desigualdades educacionais.

O uso estratégico de tecnologias educacionais para ampliar o alcance das metodologias constitui outra dimensão inovadora do projeto. A integração de AVAs, ferramentas colaborativas e plataformas de avaliação formativa ao modelo integrado de ABP e ABPr permite a criação de ambientes de aprendizagem híbridos, que combinam presencialidade e digitalidade de maneira orgânica \cite{SELWYN2022}. Essa abordagem se alinha com as tendências contemporâneas de educação digital e amplia as possibilidades de implementação em contextos com diferentes níveis de infraestrutura tecnológica.

O desenvolvimento de instrumentos de avaliação replicáveis representa uma contribuição inovadora de natureza prática. Esses instrumentos, elaborados com base em evidências científicas e validados em contextos reais, oferecem mecanismos de avaliação da efetividade das metodologias ativas integradas que podem ser adotados por diversas instituições \cite{BLACK1998}. A documentação detalhada dos processos de elaboração e de validação assegura a transparência e a reprodutibilidade dos instrumentos, facilitando sua adoção e adaptação em diferentes contextos.

A articulação entre teoria e prática, com base em evidências empíricas, constitui a quinta dimensão inovadora do projeto. A pesquisa não se limita à produção de um referencial teórico ou à geração de instrumentos práticos, mas busca articular de forma orgânica essas duas dimensões, garantindo que a teoria informe a prática e que a prática retroalimente a teoria \cite{CRESWELL2018}. Essa circularidade entre teoria e prática é um princípio fundamental das metodologias ativas e é incorporada ao próprio desenho da pesquisa de forma reflexiva.

% ----------------------------------------------------------------
\section{Identificação dos mercados}
\label{sec:mercados}

A identificação dos mercados potenciais para os resultados do projeto visa mapear os atores, as instituições e os segmentos que podem se beneficiar das contribuições teóricas e práticas desta dissertação. O reconhecimento desses mercados é essencial para a definição de estratégias de disseminação e de impacto dos resultados da pesquisa.

% ----------------------------------------------------------------
\subsection{Mercados e potencial}
\label{subsec:merc_potencial}

As instituições de ensino superior constituem um dos principais mercados-alvo para os resultados desta pesquisa. Universidades públicas e privadas em todos os estados brasileiros representam um universo amplo e diversificado de potenciais usuários dos instrumentos, protocolos e referencial teórico desenvolvidos \cite{VEIGA2014}. A crescente concorrência entre instituições de ensino superior, tanto para a atração de discentes quanto para a captação de recursos de pesquisa e de extensão, cria um incentivo para a adoção de práticas pedagógicas inovadoras que diferenciem a oferta acadêmica.

As escolas da educação básica — redes públicas e privadas de ensino fundamental e médio — constituem outro mercado significativo. A implementação da BNCC, com sua ênfase em competências e em aprendizagem significativa, cria um contexto favorável à adoção de metodologias ativas \cite{BRASIL2018}. Os protocolos e instrumentos desenvolvidos nesta pesquisa podem orientar a implementação de ABP e ABPr em escolas de diferentes portes e contextos, contribuindo para a melhoria da qualidade do ensino básico.

Órgãos públicos de educação — secretarias estaduais e municipais de educação, Ministério da Educação — constituem mercados com alto potencial de impacto. A formulação de políticas públicas de educação que incorporem metodologias ativas pode ampliar significativamente o alcance dos resultados da pesquisa, atingindo milhões de estudantes e de docentes \cite{MACHADO2025}. Entidades de formação docente — centros de formação, associações profissionais — também representam mercados importantes, na medida em que podem incorporar os referenciais teóricos e práticos desenvolvidos em seus programas de capacitação \cite{CANTANHEDE2026}.

Organizações internacionais como UNESCO, Banco Mundial e OEA constituem mercados com potencial de ampliação do impacto da pesquisa para além do contexto brasileiro. A produção de conhecimento sobre metodologias ativas em contextos de desenvolvimento pode contribuir para as agendas educacionais dessas organizações e para a disseminação de práticas inovadoras em outros países com realidades similares \cite{OECD2023}.

% ----------------------------------------------------------------
\section{Análise do setor e cadeia de valor}
\label{sec:setor_cadeia}

A análise do setor educacional brasileiro e de sua cadeia de valor busca contextualizar a proposta de pesquisa em relação à dinâmica do setor, identificando oportunidades e desafios para a implementação de metodologias ativas integradas. Essa análise é fundamentada nos referenciais estratégicos de Porter e na concepção de cadeia de valor adaptada ao contexto educacional.

% ----------------------------------------------------------------
\subsection{Estudo da posição do caso no setor, segundo o modelo das 5 forças de Porter}
\label{subsec:porter}

O setor educacional brasileiro pode ser analisado pelo modelo das cinco forças de Porter \cite{PORTER1980}, oferecendo uma compreensão estrutural das dinâmicas competitivas e das condições setoriais que influenciam a implementação de metodologias ativas. A Tabela \ref{tab:porter_cinco_forcas} sintetiza a análise das cinco forças no contexto do setor educacional brasileiro.

\begin{table}[htbp]
    \centering
    \caption{Análise das cinco forças de Porter no setor educacional brasileiro}
    \label{tab:porter_cinco_forcas}
    \small
    \begin{tabular}{p{3.5cm} p{2cm} p{8cm}}
        \toprule
        \textbf{Força} & \textbf{Intensidade} & \textbf{Caracterização no setor educacional} \\
        \midrule
        \textbf{Ameaça de novos entrantes} & Moderada & Barreiras regulatórias (MEC/CNE) protegem o setor; porém, expansão do ensino online reduz barreiras \\
        \addlinespace
        \textbf{Poder de barganha dos fornecedores} & Baixa a Moderada & Diversidade de fornecedores de tecnologia reduz poder; escassez de formadores especializados aumenta poder \\
        \addlinespace
        \textbf{Poder de barganha dos clientes} & Crescente & Acesso à informação e competição entre instituições ampliam poder de escolha dos estudantes \\
        \addlinespace
        \textbf{Ameaça de substitutos} & Alta & MOOCs, cursos online, certificações profissionais e plataformas autodidatas configuram ameaça significativa \\
        \addlinespace
        \textbf{Rivalidade entre concorrentes} & Alta & Proliferação de instituições e disputa por ranking e captação de discentes intensificam a concorrência \\
        \bottomrule
    \end{tabular}
    \par\medskip
    \footnotesize{\textit{Fonte:} Elaborado pelo autor com base em \cite{PORTER1980} e \cite{SELWYN2022}.}
    \vspace{2pt}
    \footnotesize{\textbf{Nota sobre o diálogo entre autores:} A combinação de Porter (1980) e Selwyn (2022) nesta análise é epistemologicamente deliberada. Porter fornece o modelo estrutural de análise setorial (cinco forças), enquanto Selwyn oferece a lente crítica sobre como a tecnologia transforma as dinâmicas competitivas do setor educacional. A convergência entre ambos é produtiva: Porter explica \textit{como} o setor é competitivo, e Selwyn explica \textit{por que} a tecnologia amplifica e ao mesmo tempo complexifica essa competitividade. A divergência reside na perspectiva: Porter parte de uma lógica de mercado que pode ser controversa quando aplicada a instituições públicas de educação, enquanto Selwyn questiona justamente a adequação dessa lógica ao setor educacional. A dissertação resolve essa tensão ao reconhecer que o setor educacional opera sob lógica híbrida — parcialmente competitiva (ranking, captação de discentes) e parcialmente orientada por missão social (formação cidadã, equidade).}.
\end{table}

A primeira força — a ameaça de novos entrantes — apresenta um grau moderado de intensidade no setor educacional. As barreiras regulatórias impostas pelo Conselho Nacional de Educação (CNE) e pelo Ministério da Educação (MEC), que exigem credenciamento, autorização e reconhecimento de cursos, funcionam como mecanismos de proteção significativos. Contudo, a proliferação de instituições de ensino online e a facilitação do acesso à tecnologia educacional reduzem parcialmente essas barreiras, criando um cenário de pressão competitiva crescente \cite{SELWYN2022}. No contexto específico das metodologias ativas, a entrada de novos players — como plataformas educacionais e organizações de formação docente — representa uma ameaça potencial às instituições tradicionais que não se adaptam às demandas de inovação pedagógica.

A segunda força — o poder de barganha dos fornecedores — varia significativamente de acordo com o segmento do setor. No que se refere a recursos didáticos e tecnológicos, a diversidade de opções disponíveis no mercado confere aos fornecedores um poder de barganha relativamente baixo, pois as instituições podem alternar entre múltiplos provedores de plataformas, softwares e conteúdos educacionais \cite{MEZZARI2011}. Por outro lado, no que diz respeito a formação docente especializada em metodologias ativas, o poder de barganha tende a ser maior, uma vez que profissionais qualificados para a implementação e a mediação de ABP e ABPr são relativamente escassos no mercado educacional brasileiro \cite{CANTANHEDE2026}. Essa assimetria reforça a importância da formação continuada como estratégia de redução da dependência em relação a fornecedores especializados.

A terceira força — o poder de barganha dos clientes (discentes) — apresenta uma tendência de crescimento no setor educacional brasileiro. O acesso à informação, a crescente oferta de alternativas educacionais e a competição entre instituições conferem aos estudantes uma capacidade crescente de escolha e de cobrança por qualidade \cite{MENEZES2020}. No contexto das metodologias ativas, os discentes que experimentam práticas pedagógicas inovadoras tendem a desenvolver expectativas mais elevadas em relação à qualidade da experiência de aprendizagem, o que pode funcionar como um motor de mudança para as instituições que desejam reter e atrair estudantes \cite{BACICH2018}.

A quarta força — a ameaça de produtos substitutos — é particularmente relevante no contexto educacional contemporâneo. A crescente oferta de cursos online, MOOCs (Massive Open Online Courses), plataformas de aprendizagem autodidata e certificações profissionais alternativas configura um cenário de ameaça significativa para as instituições de ensino tradicionais \cite{SELWYN2022}. No que se refere especificamente às metodologias ativas, as plataformas de aprendizagem baseada em problemas e em projetos disponíveis na internet representam potenciais substitutos parciais, especialmente para estudantes que buscam experiências de aprendizagem mais flexíveis e acessíveis. A resposta a essa ameaça passa pela diferenciação da oferta presencial e pela incorporação de elementos de metodologias ativas que não podem ser plenamente replicados em ambientes virtuais isolados.

A quinta força — a rivalidade entre concorrentes — é alta no setor educacional brasileiro, tanto no que se refere à disputa por ranking acadêmico quanto à competição por captação de discentes e por diferenciação de oferta. A proliferação de instituições de ensino superior e a ampliação do acesso à educação básica pública intensificam a concorrência em todos os segmentos do setor \cite{PORTER1980}. No contexto das metodologias ativas, a adoção de ABP e ABPr pode funcionar como um fator de diferenciação competitiva, na medida em que oferece uma experiência educacional diferenciada e de maior valor agregado para os discentes \cite{BACICH2018}. Instituições que lideram na implementação dessas metodologias podem conquistar vantagens competitivas significativas em termos de reputação, captação de estudantes e produção de conhecimento\footnote{\textbf{Diálogo filosófico — Foucault × Porter:} A análise de Porter opera dentro de uma lógica de poder institucional que Michel Foucault (1975/2014) problematizaria diretamente. Em \textit{Vigiar e Punir}, Foucault descreve como as instituições modernas — escolas, prisões, hospitais — operam por meio de ``disciplina'' — um poder que não é repressivo, mas produtivo: molda corpos, normaliza comportamentos e produz sujeitos dóceis. O modelo tradicional de ensino, centrado na exposição oral e na avaliação por provas escritas, pode ser lido como uma tecnologia disciplinar foucaultiana: o professor é o vigilante, o horário é a grade, a prova é o exame. As metodologias ativas, ao colocar o estudante no centro e ao promover a autonomia, representam uma ruptura parcial com essa lógica disciplinar — mas não uma ruptura completa, pois ainda operam dentro de estruturas institucionais que impõem grades horárias, currículos obrigatórios e avaliações padronizadas. Essa tensão entre autonomia estudantil e controle institucional é fundamental para compreender por que a implementação de metodologias ativas enfrenta resistência: não é apenas falta de formação ou infraestrutura, mas uma ameaça ao \textit{modus operandi} do poder disciplinar que sustenta o modelo tradicional. A análise de Porter explica \textit{como} o setor é competitivo; Foucault explica \textit{por que} a mudança é estruturalmente resistida.}.

% ----------------------------------------------------------------
\subsection{Estudo da cadeia de valor setorial}
\label{subsec:cadeia_valor}

A cadeia de valor do setor educacional compreende um conjunto de atividades primárias e de suporte que, articuladas de maneira coordenada, contribuem para a geração de valor educacional. As atividades primárias incluem o planejamento curricular e de ensino, que constitui a etapa fundamental de definição dos objetivos de aprendizagem, dos conteúdos e das estratégias pedagógicas \cite{BRASIL2018}. No contexto das metodologias ativas, o planejamento curricular deve contemplar a seleção de problemas autênticos e de projetos contextualizados que sirvam de eixo organizador para a aprendizagem, exigindo uma revisão profunda dos currículos tradicionais orientados por ementas disciplinares \cite{COSTA2020}.

O processo de ensino-aprendizagem constitui a atividade primária central da cadeia de valor educacional. Na perspectiva das metodologias ativas, esse processo é marcado pela centralidade do estudante, pela colaboração entre pares e pela mediação ativa do docente \cite{MORAN2018}\footnote{\textbf{Trecho original (MORAN, 2018, p. 33):} ``O professor deixa de ser o centro do processo educacional para se tornar um mediador, um facilitador da aprendizagem dos alunos.'' \textbf{Tradução:} O professor torna-se mediador e facilitador da aprendizagem. \textbf{Resenha crítica:} A citação fundamenta a transformação do papel docente no contexto de metodologias ativas. Moran é a referência mais autorizada na literatura brasileira sobre o tema. Limitação: a afirmação é conceitual, não empírica — não há dados sobre como essa transformação efetivamente ocorre na prática docente brasileira. O autor deveria complementar com estudos empíricos sobre a transição do papel docente.}. A implementação de ABP e ABPr transforma radicalmente a dinâmica da sala de aula, substituindo a exposição oral por atividades de investigação, resolução de problemas e produção de soluções. Essa transformação demanda não apenas mudanças metodológicas, mas também alterações nos espaços físicos, nos materiais didáticos e nos sistemas de avaliação \cite{BACICH2018}.

A avaliação e a certificação constituem atividades primárias que requerem adaptação significativa no contexto das metodologias ativas. A avaliação formativa, que monitora o processo de aprendizagem ao longo do tempo, deve coexistir com a avaliação somativa, que verifica a achieved de competências ao final de ciclos formativos \cite{BLACK1998}\footnote{\textbf{Diálogo entre autores:} A citação de Black e Wiliam (1998) e Soares (2013) no mesmo parágrafo estabelece um diálogo epistemológico produtivo entre a tradição anglo-saxônica e a tradição brasileira de avaliação educacional. Black e Wiliam (1998) formularam o conceito de ``assessment for learning'' — avaliação formativa como ferramenta de melhoria do desempenho — em um contexto britânico onde a avaliação era predominantemente somativa e punitiva. Soares (2013), a partir do contexto brasileiro, analisa como as práticas avaliativas nacionais (ENEM, SAEB, provas de rede) reproduzem uma lógica de classificação e ordenação que contradiz os princípios da avaliação formativa. A convergência entre ambos é clara: ambos defendem que a avaliação deve servir à aprendizagem, não à seleção. A divergência é contextual: Black e Wiliam operam em um sistema educacional com tradição de avaliação formativa no cotidiano escolar, enquanto Soares opera em um sistema onde a avaliação é centralizada e sumativa. A dissertação articula essas duas perspectivas ao propor instrumentos de avaliação que contemplem tanto a dimensão formativa (Black e Wiliam) quanto a realidade das avaliações externas brasileiras (Soares).}. No contexto de ABP e ABPr, os instrumentos de avaliação devem contemplar dimensões como a capacidade de resolução de problemas, o trabalho em equipe, a comunicação e a criatividade, além do domínio de conteúdos conceituais. Essa ampliação do escopo avaliativo representa um desafio e uma oportunidade para a melhoria da qualidade educacional \cite{SOARES2013}.

A extensão e o relacionamento com a comunidade constituem atividades primárias que ganham relevância particular no contexto das metodologias ativas. A inserção de problemas e projetos provenientes da comunidade no currículo fortalece o vínculo entre a instituição de ensino e o seu entorno social, promovendo a responsabilidade social e a pertinência da formação oferecida \cite{BARRON1998}. A gestão acadêmica e administrativa, por sua vez, deve se adaptar para suportar a implementação de metodologias ativas, contemplando aspectos como a alocação de recursos, a organização do tempo letivo e a gestão de processos que envolvam múltiplas disciplinas e atores \cite{OLIVEIRA2023}.

As atividades de suporte da cadeia de valor educacional desempenham um papel fundamental para a viabilização das metodologias ativas. A infraestrutura física e tecnológica constitui um pré-requisito essencial, na medida em que a implementação de ABP e ABPr demanda espaços flexíveis, conectividade estável e acesso a ferramentas digitais de colaboração \cite{MEZZARI2011}. A formação e o desenvolvimento docente representam outra atividade de suporte de crucial importância, uma vez que a adoção de metodologias ativas exige competências pedagógicas, tecnológicas e de mediação que frequentemente não são contempladas na formação inicial de professores \cite{CANTANHEDE2026}.

A gestão de pessoas e processos, o marketing e a comunicação institucional e as finanças e o controle orçamentário completam o conjunto de atividades de suporte da cadeia de valor educacional. A gestão de pessoas deve considerar a necessidade de formação continuada e de reconhecimento profissional para os docentes que adotam metodologias ativas, que frequentemente demandam um investimento de tempo e de esforço significativamente maior do que as práticas tradicionais \cite{MIRANDA2022}. O marketing e a comunicação institucional devem divulgar as inovações pedagógicas implementadas, atraindo estudantes e parceiros que compartilhem os valores de aprendizagem ativa e significativa \cite{PORTER1980}. As finanças e o controle orçamentário devem contemplar a alocação de recursos para infraestrutura, formação e pesquisa, assegurando a sustentabilidade financeira da implementação das metodologias ativas ao longo do tempo \cite{FREITAS2024}.

% ----------------------------------------------------------------
\subsection{Diagrama de processos do caso e contextualização de solução}
\label{subsec:diagrama_processos}

O diagrama de processos do caso contempla as etapas fundamentais para a implementação da solução proposta, articulando diagnóstico, planejamento, formação, implementação, avaliação, refinamento e disseminação. A etapa de diagnóstico envolve o levantamento de necessidades e a análise detalhada do contexto institucional, contemplando aspectos como a infraestrutura disponível, o perfil dos docentes e discentes, a cultura organizacional e as práticas pedagógicas vigentes \cite{CRESWELL2018}. Esse diagnóstico inicial é fundamental para a personalização da solução às especificidades de cada instituição e para a definição de prioridades de intervenção.

A etapa de planejamento compreende a elaboração do plano de implementação de ABP/ABPr, que deve contemplar objetivos de aprendizagem claramente definidos, seleção de problemas e projetos autênticos, cronograma de atividades, critérios de avaliação e estratégias de mediação docente \cite{COSTA2020}. O planejamento deve ser realizado de forma colaborativa, envolvendo docentes, gestores e, quando possível, discentes, promovendo a apropriação coletiva do processo e a corresponsabilidade pelos resultados.

A etapa de formação docente constitui um componente crítico para o sucesso da implementação. A capacitação dos professores deve abranger não apenas os fundamentos teóricos de ABP e ABPr, mas também aspectos práticos como o design de problemas e projetos, as estratégias de mediação, a utilização de ferramentas tecnológicas e as práticas de avaliação formativa \cite{SOARES2021}. A formação deve ser contínua e adaptativa, permitindo que os docentes aprimorem suas práticas ao longo do processo de implementação \cite{RIOS2026}.

A etapa de implementação corresponde à execução das atividades pedagógicas com metodologias ativas integradas. Nessa fase, os estudantes são desafiados a investigar problemas autênticos, a colaborar em equipes e a produzir soluções que demonstrem a aplicação dos conhecimentos adquiridos \cite{BARRON1998}. O papel do docente se transforma de transmissor de conteúdos para mediador de aprendizagem, exigindo competências de escuta, de feedback construtivo e de gestão da dinâmica de grupo \cite{MORAN2018}.

A etapa de avaliação envolve o monitoramento contínuo dos processos e dos resultados de aprendizagem, utilizando instrumentos que contemplem tanto dimensões cognitivas quanto socioemocionais e de competência \cite{BLACK1998}. A retroalimentação contínua, tanto para discentes quanto para docentes, permite identificar aspectos que requerem ajustes e aprimoramentos. A etapa de refinamento corresponde à implementação desses ajustes com base nos dados coletados, promovendo a melhoria contínua do processo de implementação \cite{CRESWELL2018}.

A etapa de disseminação consiste no compartilhamento de resultados e de boas práticas com a comunidade educacional e acadêmica. A publicação de resultados em periódicos científicos, a apresentação em eventos acadêmicos e a elaboração de relatórios técnicos para gestores educacionais constituem estratégias de disseminação que ampliam o impacto da pesquisa para além do contexto institucional onde foi conduzida \cite{SANTOS2024}. A documentação detalhada do processo de implementação facilita a replicação da experiência em outras instituições, contribuindo para a disseminação de metodologias ativas integradas no sistema educacional brasileiro.

% ============================================================
% Fim do Capítulo 2 — Aspectos Estratégicos
% ============================================================
